\notechapter{3}{Preliminary Frobenius comparison of two quantum harmonic-oscillator constructions}

\textbf{Status: proposed preliminary experiment.} The comparison, convergence
study, and plots described in this note have not yet been carried out. No
numerical values or observations are reported. The purpose of the experiment is
to compare a ladder-operator representation with a harmonic-oscillator
potential enclosed by a finite Dirichlet box after both have been placed in the
same retained ladder coordinates.

\section*{Question addressed by the experiment}

The one-dimensional quantum harmonic oscillator admits the exact
ladder-operator representation
\begin{equation}
    \hat H_{\mathrm{QHO}}
    =
    \hbar\omega
    \left(
        \hat a^\dagger\hat a+\frac{1}{2}\hat I
    \right)
\end{equation}
on $L^2(\mathbb R)$. The same potential can be placed inside a particle in a
box and represented on a finite spatial grid. Enlarging the box is expected to
reduce the influence of the Dirichlet walls on low-lying oscillator states,
while refining the grid is expected to reduce spatial-discretization error.

The central quantity is not a raw difference between a ladder matrix and a grid
matrix. Those matrices act in different coordinate spaces and may have
different dimensions. Instead, the finite-box Hamiltonian must first be pulled
back to a declared finite ladder space. The primary question is then
\begin{quote}
How does the Frobenius discrepancy between the pulled-back finite-box operator
and the exact retained ladder Hamiltonian behave as the box is enlarged and the
spatial discretization is refined?
\end{quote}

This construction also permits comparisons between spatial discretization
schemes with different grid spacings or dimensions because every represented
operator is mapped to the same retained ladder coordinates before subtraction.

\section*{Exact retained ladder Hamiltonian}

The oscillator length is
\begin{equation}
    \ell
    =
    \sqrt{\frac{\hbar}{m\omega}}.
\end{equation}
The normalized number states satisfy
\begin{align}
    \hat a\lvert n\rangle
    &=
    \sqrt{n}\,\lvert n-1\rangle,
    \\
    \hat a^\dagger\lvert n\rangle
    &=
    \sqrt{n+1}\,\lvert n+1\rangle,
\end{align}
and
\begin{equation}
    \hat H_{\mathrm{QHO}}\lvert n\rangle
    =
    E_n\lvert n\rangle,
    \qquad
    E_n
    =
    \hbar\omega
    \left(n+\frac{1}{2}\right).
\end{equation}

For a retained dimension $K$, define the abstract comparison space
\begin{equation}
    \mathcal K_K
    =
    \operatorname{span}
    \left\{
        \lvert 0\rangle,\ldots,\lvert K-1\rangle
    \right\}.
\end{equation}
In the ordered number-state basis, the exact retained Hamiltonian is
\begin{equation}
    \mathbf H_{\mathrm{lad}}^{(K)}
    =
    \hbar\omega\,
    \operatorname{diag}
    \left(
        \frac{1}{2},
        \frac{3}{2},
        \ldots,
        K-\frac{1}{2}
    \right).
\end{equation}
This is an exact spectral restriction of the real-line oscillator. It is the
reference operator for the proposed Frobenius comparison.

A finite ladder representation has an algebraic truncation boundary. If
$\mathbf a_K$ is the compressed annihilation operator, then
\begin{equation}
    [\mathbf a_K,\mathbf a_K^\dagger]
    =
    \mathbf I_K
    -
    K\lvert K-1\rangle\langle K-1\rvert.
\end{equation}
This top-state defect should be recorded when ladder operators themselves are
compared, although it does not alter the diagonal reference energies written
above.

\section*{Harmonic potential inside a Dirichlet box}

Let $X>0$ be the box half-width. The finite-box continuum Hamiltonian is
\begin{equation}
    \hat H_X
    =
    -\frac{\hbar^2}{2m}\frac{d^2}{dx^2}
    +
    \frac{1}{2}m\omega^2x^2
\end{equation}
on $L^2((-X,X))$, with domain
\begin{equation}
    \mathcal D(\hat H_X)
    =
    H^2((-X,X))\cap H_0^1((-X,X)).
\end{equation}
Thus
\begin{equation}
    \psi(-X)=\psi(X)=0.
\end{equation}
For finite $X$, this is a different continuum operator from the oscillator on
$\mathbb R$. The difference is a boundary-domain effect, not a change to the
interior multiplication potential.

For a spatial representation labeled by $h$, let the numerical state space be
$\mathcal V_{X,h}$ and let
\begin{equation}
    \mathbf H_{X,h}
    =
    \mathbf T_{\mathrm D,X,h}
    +
    \mathbf V_{\mathrm{QHO},X,h}
    \in
    \mathcal L(\mathcal V_{X,h}).
\end{equation}
Here $\mathbf T_{\mathrm D,X,h}$ represents the kinetic operator with the
declared Dirichlet closure, and $\mathbf V_{\mathrm{QHO},X,h}$ represents the
interior potential $m\omega^2x^2/2$. For a uniform grid with $N$ interior
points,
\begin{equation}
    x_j=-X+j\Delta x,
    \qquad
    \Delta x=\frac{2X}{N+1},
    \qquad
    j=1,\ldots,N.
\end{equation}
The comparison below does not require two spatial representations to have the
same $N$ or $\Delta x$.

\section*{Map from ladder states to a spatial representation}

Let
\begin{equation}
    \psi_n(x)=\langle x\mid n\rangle
\end{equation}
be the normalized real-line oscillator eigenfunction. For a uniform grid,
define the quadrature-scaled sampled vector
\begin{equation}
    \mathbf s_{n;X,h}(j)
    =
    \sqrt{\Delta x}\,\psi_n(x_j).
\end{equation}
Collect the first $K$ sampled states in
\begin{equation}
    \mathbf S_{X,h}^{(K)}
    =
    \begin{bmatrix}
        \mathbf s_{0;X,h}
        & \cdots &
        \mathbf s_{K-1;X,h}
    \end{bmatrix}.
\end{equation}
The columns need not be exactly orthonormal after restriction to a finite box
and spatial representation. Their Gram matrix is
\begin{equation}
    \mathbf G_{X,h}^{(K)}
    =
    \left(\mathbf S_{X,h}^{(K)}\right)^\dagger
    \mathbf S_{X,h}^{(K)}.
\end{equation}
Provided that $\mathbf G_{X,h}^{(K)}$ is positive definite, define the
symmetrically orthonormalized injection
\begin{equation}
    \mathbf J_{X,h}^{(K)}
    =
    \mathbf S_{X,h}^{(K)}
    \left(\mathbf G_{X,h}^{(K)}\right)^{-1/2}.
\end{equation}
It satisfies
\begin{equation}
    \left(\mathbf J_{X,h}^{(K)}\right)^\dagger
    \mathbf J_{X,h}^{(K)}
    =
    \mathbf I_K.
\end{equation}
The columns of $\mathbf J_{X,h}^{(K)}$ identify the abstract ladder coordinates
with an orthonormal $K$-dimensional subspace of $\mathcal V_{X,h}$.

The deviation
\begin{equation}
    \delta G(X,h;K)
    =
    \left\lVert
        \mathbf G_{X,h}^{(K)}-\mathbf I_K
    \right\rVert_{\mathrm F}
\end{equation}
and the condition number of $\mathbf G_{X,h}^{(K)}$ should be reported. They
measure whether finite-box restriction and spatial representation have made
the proposed ladder-state identification inaccurate or ill-conditioned. An
ill-conditioned Gram matrix is a failure of the selected comparison map, not a
small operator discrepancy.

For a nonuniform grid or another spatial scheme, the same construction requires
the declared discrete inner product or quadrature matrix. The injection must be
isometric with respect to that inner product before a Frobenius comparison is
interpreted as a Hilbert--Schmidt comparison.

\section*{Finite-box Hamiltonian in ladder coordinates}

Pull the represented finite-box Hamiltonian back to the retained ladder space:
\begin{equation}
    \mathbf A_{X,h}^{(K)}
    =
    \left(\mathbf J_{X,h}^{(K)}\right)^\dagger
    \mathbf H_{X,h}
    \mathbf J_{X,h}^{(K)}.
\end{equation}
Both
\begin{equation}
    \mathbf A_{X,h}^{(K)},
    \mathbf H_{\mathrm{lad}}^{(K)}
    \in
    \mathbb C^{K\times K}
\end{equation}
now represent operators in the same declared ladder coordinates. Their
subtraction is therefore defined.

The primary absolute discrepancy is
\begin{equation}
    D_F(X,h;K)
    =
    \left\lVert
        \mathbf A_{X,h}^{(K)}
        -
        \mathbf H_{\mathrm{lad}}^{(K)}
    \right\rVert_{\mathrm F}.
\end{equation}
A dimensionless relative discrepancy is
\begin{equation}
    \delta_F(X,h;K)
    =
    \frac{
        \left\lVert
            \mathbf A_{X,h}^{(K)}
            -
            \mathbf H_{\mathrm{lad}}^{(K)}
        \right\rVert_{\mathrm F}
    }{
        \left\lVert
            \mathbf H_{\mathrm{lad}}^{(K)}
        \right\rVert_{\mathrm F}
    }.
\end{equation}
This is the principal proposed output of the experiment. Its interpretation is
the finite-box-plus-representation discrepancy on the declared $K$-state
ladder subspace. It is not a norm difference between the full unbounded
continuum Hamiltonians.

Because the reference matrix is diagonal, the discrepancy separates as
\begin{align}
    D_F(X,h;K)^2
    &=
    D_{\mathrm{diag}}(X,h;K)^2
    +
    D_{\mathrm{off}}(X,h;K)^2,
    \\
    D_{\mathrm{diag}}(X,h;K)^2
    &=
    \sum_{n=0}^{K-1}
    \left|
        A_{X,h;nn}^{(K)}-E_n
    \right|^2,
    \\
    D_{\mathrm{off}}(X,h;K)^2
    &=
    \sum_{\substack{m,n=0\\m\neq n}}^{K-1}
    \left|
        A_{X,h;mn}^{(K)}
    \right|^2.
\end{align}
The diagonal term measures errors in the ladder-coordinate energy entries. The
off-diagonal term measures mixing produced by the finite boundary, spatial
representation, and finite-sample identification. Reporting only eigenvalue
errors would omit this off-diagonal operator information.

\section*{Comparing different spatial discretizations}

Let $h$ and $h'$ denote two spatial representations. Their grid matrices may
have different dimensions, so the raw expression
\begin{equation}
    \mathbf H_{X,h}-\mathbf H_{X,h'}
\end{equation}
is generally undefined. Each operator is instead pulled back independently:
\begin{align}
    \mathbf A_{X,h}^{(K)}
    &=
    \left(\mathbf J_{X,h}^{(K)}\right)^\dagger
    \mathbf H_{X,h}
    \mathbf J_{X,h}^{(K)},
    \\
    \mathbf A_{X,h'}^{(K)}
    &=
    \left(\mathbf J_{X,h'}^{(K)}\right)^\dagger
    \mathbf H_{X,h'}
    \mathbf J_{X,h'}^{(K)}.
\end{align}
The cross-discretization discrepancy is then
\begin{equation}
    D_F(h,h';X,K)
    =
    \left\lVert
        \mathbf A_{X,h}^{(K)}
        -
        \mathbf A_{X,h'}^{(K)}
    \right\rVert_{\mathrm F}.
\end{equation}
This remains well-defined when the two grids have different spacings, numbers
of points, stencils, or other declared representation details. The ladder
space supplies the common coordinates.

The maps $\mathbf J_{X,h}^{(K)}$ and $\mathbf J_{X,h'}^{(K)}$ are part of the
comparison and must be retained with the result. Equal output dimensions do not
make the comparison independent of these maps.

\section*{Convergence questions}

Use the dimensionless controls
\begin{equation}
    b=\frac{X}{\ell},
    \qquad
    \eta=\frac{\Delta x}{\ell}
\end{equation}
for a uniform grid. For fixed $K$, the expected limiting behavior is
\begin{equation}
    \delta_F(X,h;K)
    \longrightarrow
    0
\end{equation}
when the box expands and the spatial representation is refined in a controlled
manner. This is expected behavior to be tested, not a result established by
this note.

The two limiting effects must be separated:
\begin{enumerate}
    \item \textbf{Spatial refinement at fixed box.} Hold $b$ fixed and reduce
          $\eta$ to estimate representation error for the finite-box operator.
    \item \textbf{Box expansion at controlled resolution.} Increase $b$ while
          keeping $\eta$ fixed or sufficiently refined to estimate the effect
          of the Dirichlet walls.
    \item \textbf{Retained-space sensitivity.} Repeat the two studies for
          several values of $K$. Higher oscillator states extend farther in
          space and are expected to require larger boxes and finer resolution.
\end{enumerate}
Holding $N$ fixed while increasing $X$ does not isolate box convergence because
$\Delta x=2X/(N+1)$ then grows and the spatial resolution becomes worse.

A nested study may first estimate the refined-grid limit at each fixed $X$ and
then examine the large-box behavior of those estimates. The reverse order may
also be reported, but the order of limits and the numerical controls must be
explicit.

\section*{Plot placeholders}

\begin{enumerate}
    \item \textbf{Placeholder---primary Frobenius convergence.} Plot
          $\delta_F(X,h;K)$ against $b$ for several controlled values of
          $\eta$ and $K$.
    \item \textbf{Placeholder---fixed-box spatial refinement.} Plot
          $\delta_F(X,h;K)$ against $\eta$ at fixed $b$ for several retained
          dimensions.
    \item \textbf{Placeholder---two-parameter convergence surface.} Display
          $\delta_F$ over the $(b,\eta)$ plane at fixed $K$.
    \item \textbf{Placeholder---diagonal and off-diagonal decomposition.}
          Plot $D_{\mathrm{diag}}$ and $D_{\mathrm{off}}$ separately against
          $b$ and $\eta$.
    \item \textbf{Placeholder---operator-difference heat map.} Display the
          entries of
          $\left(\mathbf A_{X,h}^{(K)}-
          \mathbf H_{\mathrm{lad}}^{(K)}\right)/(\hbar\omega)$ for selected
          parameter values.
    \item \textbf{Placeholder---cross-grid comparison.} Plot
          $D_F(h,h';X,K)/(\hbar\omega)$ for pairs of spatial representations
          after both have been pulled back to the same ladder space.
    \item \textbf{Placeholder---map quality.} Plot $\delta G(X,h;K)$ and the
          condition number of $\mathbf G_{X,h}^{(K)}$ against $b$, $\eta$, and
          $K$.
    \item \textbf{Placeholder---retained-dimension sensitivity.} Plot the
          normalized discrepancy against $K$ at selected converged and
          unconverged box and grid settings.
\end{enumerate}
Each completed plot should identify $b$, $\eta$, $K$, the spatial inner
product, the orthonormalization convention, and the status of the displayed
values. A plot placeholder is not a calculated result.

\section*{Interpretation boundary}

The proposed experiment tests whether represented finite-box harmonic
oscillators approach the exact retained ladder operator under an explicit
state-space map. It separates boundary-domain effects, spatial-representation
error, retained-space sensitivity, and map conditioning. The Frobenius norm is
appropriate only after the two operators have been placed in the same
orthonormal ladder coordinates, as described in
Appendix~\ref{app:hilbert-schmidt-and-frobenius}.

Successful convergence would constitute numerical verification for this
illustrative oscillator comparison under the declared conditions. It would not
validate a semiconductor model, establish a continuum limit for silicon, or
supply evidence for an unperformed first-principles calculation.
